\chapter{Теоретическая часть}

\section{Определение $\gamma$~-~доверительного интервала для значения параметра распределения случайной величины}

Пусть:
1) $X$ -- случайная величина, закон распределения которой
известен с точностью до неизвестного параметра $\theta$.

Закон распределения с.в. $X$ известен с точностью до $\theta$
означает, что известен общий вид закона распределения с.в. $X$, но этот закон
зависит от неизвестного параметра $\theta$.
Если задать некоторое значение $\theta$, 
то закон распределения с.в. X будет известен полностью

Интервальной оценкой с уровнем доверия $\gamma$ 
($\gamma$~-~доверительной интервальной оценкой) 
параметра $\theta$ называют пару статистик 
$\underline{\theta}(\vec X), \overline{\theta}(\vec X)$ таких, что

\begin{equation*}
    P\{\underline{\theta}(\vec X)<\theta<\overline{\theta}(\vec X)\}=\gamma
\end{equation*}

% Поскольку границы интервала являются случайными величинами, то для различных реализаций случайной выборки $\vec X$ статистики $\underline{\theta}(\vec X), \overline{\theta}(\vec X)$ могут принимать различные значения.

Доверительным интервалом с уровнем доверия $\gamma$ ($\gamma$-доверительным интервалом) называют интервал $(\underline{\theta}(\vec x), \overline{\theta}(\vec x))$, отвечающий выборочным значениям статистик $\underline{\theta}(\vec X), \overline{\theta}(\vec X)$.

\section{Формулы для вычисления границ \\ $\gamma$-доверительного интервала для математического ожидания и дисперсии нормальной случайной величины}

Статистику $\hat{\theta}(\vec{X})$ называют точечной оценкой параметра $\theta$,
если ее выборочное значение принимается в качестве параметра $\theta$.
Т.е. $\theta = \hat{\theta}(\vec{x})$ 

Формулы для вычисления границ $\gamma$-доверительного интервала для математического ожидания при известной дисперсии:

\begin{equation}
	\label{eq:low_mu}
	\underline{\mu}(\vec{X}_{n}) = \overline{X}-\frac{S(\vec{X}_{n})}{\sqrt{n}}t_{1-\alpha}(n-1)
\end{equation}

\begin{equation}
	\label{eq:high_mu}
	\overline{\mu}(\vec{X}_{n}) = \overline{X}+\frac{S(\vec{X}_{n})}{\sqrt{n}}t_{\alpha}(n-1),
\end{equation}

\noindent где:
\begin{itemize}
	\item $\overline{X}$ -- среднее значение выборки;
	\item $n$ -- число опытов;
	\item $S(\vec{X}_{n})$ -- точечная оценка дисперсии случайной выборки $\vec{X}_{n}$;
	\item $t_{\alpha}(n-1)$ квантиль уровня $\alpha$ для распределения Стьюдента с $n-1$ степенями свободы;
	\item $\alpha$ -- величина, равная $\displaystyle\frac{(1-\gamma)}{2}$.
\end{itemize}

Формулы для вычисления границ $\gamma$-доверительного интервала для дисперсии:

\begin{equation}
	\label{eq:low_sigma}
	\underline{\sigma^{2}}(\vec{X}_{n})  =\frac{S(\vec{X}_{n})(n-1)}{h_{1-\alpha}^{n-1}}
\end{equation}

\begin{equation}
	\label{eq:high_sigma}
	\overline{\sigma^{2}}(\vec{X}_{n})  =\frac{S(\vec{X}_{n})(n-1)}{h_{\alpha}^{n-1}},
\end{equation}

\noindent где:
\begin{itemize}
	\item $n$ -- объем выборки;
	\item $h_{\alpha}^{n-1}$ -- квантиль уровня $\alpha$ распределения $\chi^{2}$ с $n-1$ степенями свободы;
	\item $\alpha$ -- величина, равная $\displaystyle\frac{(1-\gamma)}{2}$.
\end{itemize}